\documentclass{article}
\usepackage[paperwidth=180mm,paperheight=180mm,margin=16mm]{geometry}
\usepackage{amsmath,unicode-math,xcolor}
\setmathfont{XITS Math}
\usepackage{experimental/flat-core}
\pagestyle{empty}
\newcommand\UnequalFormula{%
 \begin{flatmatchedsplit}
  F&=a+b\\
   &+\frac{\displaystyle\sum_{k=1}^{n}x_k}{\displaystyle 1+\frac{p}{q}}\\
   &+z_i^2
 \end{flatmatchedsplit}}
\begin{document}
\noindent\textbf{LuaLaTeX: three rows with different heights}\par\medskip
\noindent\textcolor{blue}{Blue: actual row boxes.}\par
\noindent\textcolor{green!50!black}{Green: enclosing formula box.}
\par\noindent\textcolor{red!80!black}{Red dashed line: \texttt{\&} column boundary; crosses: row baselines.}
\[
 \begingroup\fboxsep=3pt\fboxrule=.3pt
 \fcolorbox{green!50!black}{white}{$\displaystyle\FlatShowBoxes\FlatShowAlignment\UnequalFormula$}
 \endgroup
\]
\noindent\textbf{The same output without outlines}
\[\UnequalFormula\]
\noindent Row 1 contains ordinary letters. Row 2 contains a sum with limits
inside a nested fraction. Row 3 contains a superscript and a subscript.
\par\medskip
\noindent The prototype uses each row's own height and depth. The baseline
distance is the larger of 15pt and the previous row's depth plus the next
row's height plus 3pt. It is not a uniform baseline grid.
\end{document}
